This is the related work chapter.

\subsection{Software Engineering in Education}


When people think of software engineering, they often focus on programming, but software engineering education prepares students not only to write code but also to design, develop, and maintain complex software systems.

\subsubsection{Programming Education}


Though programing is not the main focus of software engineering it serves as the entry point to software engineering, where students learn fundamental concepts that later extend into more advanced topics such as software design and system architecture.


Because programming serves as the entry point to software engineering, the choice of programming language is important, as languages with a steep learning curve may discourage beginners from engaging with the field.


The specific language to use for this entry point has been a history of debate\cite{villagomez-palacios_bridging_2025}, as language preference is subject to personal preference and efficiency varies between learners.
George Milbrandt\cite{mannila_objective_2006} suggested these criteria for language selection for beginners: it should be easy to use, have a structured design with one-entry/one-exit flow, offer powerful computing capacity while maintaining simple syntax, support variable declaration and expressive, meaningful variable names, provide easy input/output and output formatting, give immediate feedback, and include good diagnostic tools for testing and debugging.
Although these are good criteria for a language to have its leaves very little tooling for comparing different languages and their capacity to teach.
This is something that Mannila and Raadt have sought to rectify in their paper \textit{ "An Objective Comparison of Languages for  Teaching Introductory Programming"}\cite{mannila_objective_2006}, wherein they compare different programming language to find the one which is the best for teaching beginners.
They have collected 17 different criteria drawn from the design decisions made by creators of languages that are considered ``teaching languages''.




\subsubsection{Challenges in Learning Software Engineering}


\subsubsection{Computational Thinking and Conceptual Learning}


\subsubsection{Game-Based Learning}